% ch7.tex — 练习与补充 III（中文版）
\chapter{练习与补充 III}

\section{Hermite 方程的第二个解}

Hermite 方程
\begin{equation}
y'' - 2xy' + 2ny = 0
\label{eq:7.1.1}
\end{equation}
对整数 $n$ 有多项式解 $H_n(x)$。第二个线性无关解（非多项式）为
\begin{equation}
y_2(x) = H_n(x)\int^x \frac{e^{t^2}}{H_n(t)^2}\,dt,
\label{eq:7.1.2}
\end{equation}
由降阶法得到。对 $n=0$，
$y_2(x)=\erfi(x)$（虚误差函数）。

\begin{exercise}
验证当 $y_1(x)=H_n(x)$ 时，\eqref{eq:7.1.2} 中的 $y_2(x)$ 满足 Hermite 方程。
\end{exercise}

\section{高阶 Hermite 多项式}

广义 Hermite 多项式 $H_n^{(m)}(x)$ 由生成函数
\begin{equation}
\sum_{n=0}^{\infty} H_n^{(m)}(x)\,\frac{t^n}{n!}
= \exp\!\left(-t^m + m\,t\,x\right),
\label{eq:7.2.1}
\end{equation}
定义，并满足高阶热方程
\begin{equation}
\frac{\partial F}{\partial t} = \frac{\partial^m F}{\partial x^m},
\qquad F(x,0)=x^n.
\label{eq:7.2.2}
\end{equation}

其显式形式为：
\begin{equation}
H_n^{(m)}(x) = n!\sum_{r=0}^{\lfloor n/m\rfloor}
\frac{(-1)^r x^{n-mr}}{(n-mr)!\,r!}.
\label{eq:7.2.3}
\end{equation}

\begin{exercise}
证明 $H_n^{(3)}(x)$ 满足 $\partial_t F = \partial_x^3 F$，并推导联系 $H_{n+1}^{(3)}$、$H_n^{(3)}$ 与 $H_{n-2}^{(3)}$ 的递推关系。
\end{exercise}

\section{多指标 Hermite 多项式}

对于 $d$ 维矢量 $\vect x=(x_1,\dots,x_d)$ 与多指标
$\vect n=(n_1,\dots,n_d)$，
\begin{equation}
H_{\vect n}(\vect x) = (-1)^{|\vect n|}
e^{x_1^2+\cdots+x_d^2}
\frac{\partial^{|\vect n|}}{\partial x_1^{n_1}\cdots\partial x_d^{n_d}}
e^{-(x_1^2+\cdots+x_d^2)},
\label{eq:7.3.1}
\end{equation}
其中 $|\vect n|=n_1+\cdots+n_d$。

生成函数：
\begin{equation}
\sum_{\vect n}\frac{\vect t^{\vect n}}{\vect n!}\,
H_{\vect n}(\vect x)
= \exp\!\left(\vect t\cdot\vect x - \tfrac12\vect t\cdot\vect t\right),
\label{eq:7.3.2}
\end{equation}
其中 $\vect n! = n_1!\cdots n_d!$，$\vect t^{\vect n}=t_1^{n_1}\cdots t_d^{n_d}$。

\begin{exercise}
证明多维 Rodrigues 公式 \eqref{eq:7.3.1}，并推导权重为 $e^{-|\vect x|^2}$ 的正交关系。
\end{exercise}

\section{产生--湮灭算子代数及其物理应用}

$d$ 维谐振子的升降算子：
\begin{equation}
\hat a_i = \frac{1}{\sqrt{2}}\!\left(x_i + \frac{\partial}{\partial x_i}\right),\qquad
\hat a_i^\dagger = \frac{1}{\sqrt{2}}\!\left(x_i - \frac{\partial}{\partial x_i}\right),
\label{eq:7.4.1}
\end{equation}
且 $[\hat a_i,\hat a_j^\dagger]=\delta_{ij}$，$[\hat a_i,\hat a_j]=0$。

哈密顿量因式分解：
\begin{equation}
\hat H = \frac12\sum_{i=1}^d (\hat a_i\hat a_i^\dagger + \hat a_i^\dagger\hat a_i)
= \sum_{i=1}^d \left(\hat a_i^\dagger\hat a_i + \frac12\right).
\label{eq:7.4.2}
\end{equation}

$SU(1,1)$ 实现：
$\hat K_+=\frac12\hat a^\dagger\hat a^\dagger$，
$\hat K_-=\frac12\hat a\hat a$，
$\hat K_0=\frac14(\hat a^\dagger\hat a^\dagger\hat a\hat a + \hat a\hat a\hat a^\dagger\hat a^\dagger)$，
满足 $[\hat K_0,\hat K_\pm]=\pm\hat K_\pm$，
$[\hat K_+,\hat K_-]=-2\hat K_0$。

压缩算子：
$\hat S(\xi)=\exp(\xi\hat K_+ - \xi^*\hat K_-)$，
它把 $\hat a$ 变换为 $\cosh r\,\hat a + e^{i\theta}\sinh r\,\hat a^\dagger$。

\begin{exercise}
证明 BCH 公式给出
$e^{\xi\hat K_+-\xi^*\hat K_-}=e^{\tau\hat K_+}\,e^{\ln(1-|\tau|^2)\hat K_0}\,e^{-\tau^*\hat K_-}$，
其中 $\tau = \tanh r\,e^{i\theta}$。
\end{exercise}

\section{Eisenstein 整数}

环 $\mathbb Z[\omega]$（$\omega=e^{2\pi i/3}$，Eisenstein 整数）是 umbral 移位演算的基础。元素形如 $a+b\omega$，其中
$a,b\in\mathbb Z$ 且 $\omega^2+\omega+1=0$。

范数 $N(a+b\omega)=a^2-ab+b^2$ 是可乘的。
单位元为 $\{\pm1,\pm\omega,\pm\omega^2\}$。

在 umbral 演算中，满足 $\hat\theta x^n = n x^{n-1}$ 的移位算子 $\hat\theta$ 可以形式地与 Eisenstein 整数等同，给出分数导数的离散模型：
\begin{equation}
\hat\theta^\alpha x^n = \frac{\Gamma(n+1)}{\Gamma(n-\alpha+1)}\,x^{n-\alpha},
\qquad \alpha\in\mathbb Z[\omega].
\label{eq:7.5.1}
\end{equation}

\begin{exercise}
验证 $\mathbb Z[\omega]$ 中的乘法规则，并证明范数可乘：$N(z_1z_2)=N(z_1)N(z_2)$。
\end{exercise}

\section{谐振子哈密顿量的形式方面与进一步的杂项考虑}

通过因式分解的代数求解：
\begin{equation}
\hat H = \frac12(\hat a\hat a^\dagger + \hat a^\dagger\hat a)
= \hat a^\dagger\hat a + \frac12,
\qquad \hat a = \frac{1}{\sqrt{2}}\!\left(x+\frac{d}{dx}\right).
\label{eq:7.6.1}
\end{equation}

超对称伙伴势：
$\hat H_-=\hat a^\dagger\hat a$，$\hat H_+=\hat a\hat a^\dagger$，
且 $\hat H_+\psi_n^{(+)}=\psi_{n+1}^{(-)}$ 等等。

相干态 $|\alpha\rangle=e^{-|\alpha|^2/2}\sum_{n=0}^\infty\frac{\alpha^n}{\sqrt{n!}}|n\rangle$
满足 $\hat a|\alpha\rangle=\alpha|\alpha\rangle$，是湮灭算子的本征态。

\begin{exercise}
证明相干态不正交：
$\langle\beta|\alpha\rangle = \exp(-\tfrac12|\alpha|^2-\tfrac12|\beta|^2+\beta^*\alpha)$。
\end{exercise}

\section{含时哈密顿量}

对于 $\hat H(t)=\hbar\vect\Omega(t)\cdot\vect\sigma$，演化算子是一个时间排序指数
\begin{equation}
\hat U(t) = \mathcal T\exp\!\left(-\frac{i}{\hbar}\int_0^t\hat H(t')\,dt'\right),
\label{eq:7.7.1}
\end{equation}
用第 2 章的 Wei--Norman 方法解缠。

对于哈密顿量 $\hat H(t)=\omega(t)\hat\sigma_z+\Omega(t)\hat\sigma_x$，
演化算子可以写成
$\hat U(t)=e^{i\alpha(t)\hat\sigma_z}e^{\beta(t)\hat\sigma_x}e^{i\gamma(t)\hat\sigma_z}$，
其中 $\alpha,\beta,\gamma$ 通过求解运动方程确定。

\section{Dyson 级数}

演化算子的 Dyson 展开为
\begin{equation}
\hat U(t)=\hat 1+\sum_{n=1}^{\infty}(-i)^n\int_0^t dt_1\cdots\int_0^{t_{n-1}}dt_n\,
\hat H(t_1)\cdots\hat H(t_n).
\label{eq:7.8.1}
\end{equation}
这是 $i\partial_t\hat U=\hat H\hat U$（$\hat U(0)=\hat 1$）的微扰解。

\subsection{Dyson 展开之外}

Magnus 展开：
\begin{equation}
\hat U(t) = \exp\!\left(\sum_{n=1}^{\infty}\hat\Omega_n(t)\right),
\label{eq:7.8.2}
\end{equation}
其中 $\hat\Omega_1=-i\int_0^t\hat H(t_1)dt_1$，
$\hat\Omega_2=-\frac12\int_0^t dt_1\int_0^{t_1}dt_2\,[\hat H(t_1),\hat H(t_2)]$，等等。

Wei--Norman 方法：写 $\hat U(t)=\prod_j e^{f_j(t)\hat G_j}$，其中
$\{\hat G_j\}$ 是一个 Lie 代数基；函数 $f_j(t)$ 满足一个封闭的一阶 ODE 组。

BCH 公式：
\begin{equation}
e^{\hat A}e^{\hat B} = \exp\!\left(\hat A+\hat B
+ \frac12[\hat A,\hat B]
+ \frac{1}{12}[\hat A,[\hat A,\hat B]]
+ \frac{1}{12}[\hat B,[\hat B,\hat A]]+\cdots\right).
\label{eq:7.8.3}
\end{equation}

\begin{exercise}
通过要求 $\hat U(t)=\exp(\hat\Omega_1+\hat\Omega_2+\cdots)$ 并按 $t$ 逐阶匹配，从 Dyson 级数推导 Magnus 展开的前两项。
\end{exercise}

\section{特殊多项式与微扰论}

Hermite/Laguerre 多项式生成受微扰谐振子的矩阵元，把代数方法与 Rayleigh--Schr\"odinger 微扰论联系起来。

对于 $\hat H = \hat H_0 + \lambda\hat V$，能量的第 $n$ 阶修正用矩阵元
$\langle m|\hat V|n\rangle$ 表示，对谐振子本征态而言这些矩阵元是 $\lambda$ 的多项式，其系数涉及 Hermite 多项式。

微扰级数
$E_n(\lambda)=\sum_{k=0}^\infty E_n^{(k)}\lambda^k$
的系数为
\begin{equation}
E_n^{(k)} = \langle n|\hat V|n\rangle^{(k)},
\label{eq:7.9.1}
\end{equation}
其中第 $k$ 阶项由更低阶的波函数修正构造。

\begin{exercise}
对于 $\hat H_0=\frac12(\hat p^2+\hat x^2)$ 与 $\hat V=\hat x^4$，
用 Hermite 多项式矩阵元计算 $E_0^{(1)}$ 与 $E_0^{(2)}$。
\end{exercise}

% 图形
\begin{figure}[h]
\centering
\includegraphics[max width=0.85\linewidth,keepaspectratio]{fig7_1.png}
\caption{产生--湮灭算子代数：能级图。}
\label{fig:7.1}
\end{figure}
